% Fichier complété par tools/complete_missing_corrections.py.
% Source : DYN/DYN-04-TorseurDynamique/50_BancBalafre/50_BancBalafre.tex
% Statut : rédigé (5 question(s)).
\begin{corrigebox}[Corrigé rédigé]
\CorrigeQuestion{1}
Le coeur de butée effectue une translation suivant l'axe de guidage;
                son orientation reste constante par rapport au bâti.
\CorrigeQuestion{2}
La géométrie du guidage impose une relation linéaire
                $v(t)=k_y\dot y(t)$, où $k_y$ est le rapport cinématique lu sur le schéma
                (égal à 1 si $y$ est directement la coordonnée du centre).
\CorrigeQuestion{3}
En $G_{CB}$ :
                $\{\mathcal D(CB/0)\}=\{M_{CB}\dot v\vec e;
                \vec0\}_{G_{CB}}$ pour une translation rectiligne.
\CorrigeQuestion{4}
En $G_{JR}$, la partie mobile suit la même translation :
                $\{\mathcal D(JR/0)\}=\{M_{JR}\dot v\vec e;\vec0\}_{G_{JR}}$.
\CorrigeQuestion{5}
Après transport au point $G$ et sommation,
                $\{\mathcal D(S/0)\}_G=\{(M_{CB}+M_{JR})\dot v\vec e;\vec0\}_G$
                lorsque $G$ est le centre d'inertie de l'ensemble et qu'il n'y a pas de
                rotation.
\end{corrigebox}
